\subsection{Proof of Proposition~\ref{prop:nearcrit}}\label{app:nearcrit}

Write $D_X(z)=z-a_X-\kappa G_X(z)$ for economy $X$, so that $H_X=\kappa z/D_X$ and the stationary price spectrum is $f_X(\omega)=\sigma^2\kappa^2/|D_X(e^{i\omega})|^2$. Because both economies have budget $1-\epsilon$, $D_A(1)=D_B(1)=\epsilon$.

\emph{Step 1: the difference vanishes at zero frequency.} $D_B-D_A=-(a_B-a_A)-\kappa(G_B-G_A)$, and the choice of persistences gives $a_B-a_A=-\kappa(G_B(1)-G_A(1))$. With $E=G_B-G_A$ this is $\Delta(z)\equiv D_B(z)-D_A(z)=-\kappa\,(E(z)-E(1))$, which does not depend on $\epsilon$. Write $E(z)=\sum_me_mz^{-m}$ with $e_m=k^B_m-k^A_m$. Then $|E(e^{i\omega})-E(1)|\le|\omega|\sum_mm|e_m|$. The rates are $r_c+\Delta\eta\,s$, and expanding $(r_c+\Delta\eta\,s)^m$ in powers of $\Delta\eta$, the moments of order up to $2K-2$ of the two normalized rules cancel, so $|e_m|\le2W\sum_{j\ge2K-1}\binom mj r_c^{m-j}(\Delta\eta)^j$. Summing, $\sum_mm|e_m|\le C_1(\Delta\eta)^{2K-1}$, with $C_1$ depending on $K,W,r_c$ and a bound on $\Delta\eta$. Hence $|\Delta(e^{i\omega})|\le\kappa C_1(\Delta\eta)^{2K-1}|\omega|$.

\emph{Step 2: the denominator is bounded below by $\epsilon+|\omega|$.} Let $a^*=\lim_{\epsilon\to0}a_A(\epsilon)$ and $D^*=D_A-\epsilon$, the denominator of economy $A$ at budget one; $D_A=D^*+\epsilon$. At budget one the closed-loop matrix $J$ of the price and the groups is entrywise nonnegative with a positive diagonal and positive first row and column (since $a^*>0$ and all weights are positive), so it is primitive. Its Perron root is $1$, it is simple, and no other eigenvalue lies on the unit circle. The zeros of $D^*$ in $|z|\ge1$ are eigenvalues of $J$, so $D^*$ has no zero on the unit circle except $z=1$, and there $D^{*\prime}(1)=1+\kappa\sum_iw_ir_i/(1-r_i)^2>0$. Hence there are $\delta,c_1,C_2,m>0$ with $|\mathrm{Im}\,D^*(e^{i\omega})|\ge c_1|\omega|$ and $\mathrm{Re}\,D^*(e^{i\omega})\ge-C_2\omega^2$ for $|\omega|\le\delta$, and $|D^*(e^{i\omega})|\ge m$ for $\delta\le|\omega|\le\pi$. For $|\omega|\le\delta$: if $\epsilon\ge2C_2\omega^2$ then $\mathrm{Re}\,D_A\ge\epsilon/2$ and $|D_A|\ge\max(\epsilon/2,c_1|\omega|)$; if $\epsilon<2C_2\omega^2$ then $|D_A|\ge c_1|\omega|$ and $\epsilon<2C_2\delta|\omega|$. In both cases $|D_A(e^{i\omega})|\ge c_2(\epsilon+|\omega|)$. For $|\omega|\ge\delta$ and $\epsilon\le m/2$, $|D_A|\ge m/2$. So $|D_A(e^{i\omega})|\ge c_3(\epsilon+|\omega|)$ on $[-\pi,\pi]$ for all $\epsilon\le\epsilon_1$. These constants can be taken uniform in $\Delta\eta\in(0,\Delta\eta_0]$. As $\Delta\eta\to0$ the kernel of $A$ converges to that of the one-group economy at rate $r_c$ in the norm $\sum_m(1+m^2)|k_m|$, so $D^*$ and its first two derivatives converge uniformly on the circle. The limit has the same properties: its closed-loop matrix is primitive with Perron root $1$, and its derivative at $1$ is at least $1$. Hence $c_1$, $C_2$, $\delta$ and $m$ can be chosen for the limit and kept, after halving $c_1$ and $m$ and doubling $C_2$, for all small $\Delta\eta$.

\emph{Step 3: the bound.} By Steps 1 and 2, $\rho(\omega)=|\Delta/D_A|\le\kappa C_1(\Delta\eta)^{2K-1}|\omega|/(c_3(\epsilon+|\omega|))\le C_3(\Delta\eta)^{2K-1}$, uniformly in $\epsilon$ and $\omega$. The ratio of spectra is $u=f_B/f_A=|D_A/D_B|^2=|1+\Delta/D_A|^{-2}$, so for $\rho\le1/8$, $|u-1|\le4\rho$. The divergence rate of stationary Gaussian processes is $D_\epsilon(B\|A)=\frac1{4\pi}\int_{-\pi}^{\pi}(u-1-\log u)\,d\omega$, and $u-1-\log u\le2(u-1)^2$ for $|u-1|\le1/2$. Hence $D_\epsilon(B\|A)\le16C_3^2(\Delta\eta)^{4K-2}$ for all $\epsilon\le\min(\epsilon_0,\epsilon_1)$ and all $\Delta\eta$ small enough that $C_3(\Delta\eta)^{2K-1}\le1/8$. For $\epsilon_1<\epsilon\le\epsilon_0$ both economies are uniformly stable, so $|D_A|$ is bounded below by a constant on the circle, while $|\Delta|\le|a_B-a_A|+\kappa|E(e^{i\omega})|=O((\Delta\eta)^{2K-1})$; Step~3 then applies with this lower bound in place of Step~2. Steps 2 and 3 are machine-checked, with Whittle's formula, Step 1 and the local properties of $D^*$ as hypotheses (\texttt{p1\_nearcrit\_lower}, \texttt{p1\_ratio\_bound}, \texttt{p1\_nearcrit\_kl\_bound}, \nolinkurl{Paper1NearCritical.lean}). \qed

\paragraph{Numerical check.} For $K=2,3$, $\Delta\eta\in\{0.1,0.05,0.025\}$ and margins from $10^{-1}$ to $10^{-6}$, the ratio $D_\epsilon(B\|A)/(\Delta\eta)^{4K-2}$ increases to a finite limit as $\epsilon\to0$ (0.41 to 0.48 for $K=2$, 0.65 to 0.71 for $K=3$; \nolinkurl{fragility_uniform_check.py}). With a common persistence instead of a common budget, the divergence rate grows like $\epsilon^{-1}$, which is the budget difference becoming visible. This proposition is not machine-checked.


\paragraph{Finite samples (numerical).} The proposition bounds the information per observation. A finite-sample bound uniform in $\epsilon$ is not proved. Numerically, the exact Kullback--Leibler divergence between the laws of $n$ consecutive stationary observations, computed from the state-space autocovariances, lies between $97.8\%$ and $99.8\%$ of $n$ times the per-observation rate in every case computed: $K=2$, $\Delta\eta\in\{0.1,0.2\}$, $\epsilon\in\{10^{-1},10^{-2},10^{-3}\}$ and $n\in\{250,1000,4000\}$, at the parameters of Figure~\ref{fig:barrier} (\nolinkurl{research_extension/fragility_reveals/finite_sample_check.py}). So, in these cases, data requirements computed from the per-observation rate are not loosened by finite samples.
